% Part: computability
% Chapter: recursive-functions
% Section: general-recursive-functions

\documentclass[../../../include/open-logic-section]{subfiles}

\begin{document}

\olfileid{cmp}{rec}{gen}
\olsection{一般递归函数}

获取全函数集合还有另一种方法。如果一个全函数 $f(x,\vec z)$ 对每个自然数序列 $\vec z$，都存在 $x$ 使得 $f(x,\vec z) = 0$，则称它是\emph{正则的}。换言之，正则函数恰好就是那些对其应用无界搜索后，最终仍能得到全函数的函数。因此，我们可以保守地将无界搜索限制于正则函数：

\begin{defn}
\ollabel{defn:general-recursive}
\emph{一般递归函数}集合是满足以下条件的最小函数集合（各元数的从自然数到自然数的函数）：包含零函数、后继函数与投影函数，且在复合、原始递归以及对\emph{正则}函数的无界搜索下封闭。
\end{defn}

显然，每个一般递归函数都是全函数。\olref{defn:general-recursive} 与 \olref[par]{defn:recursive-fn} 的区别在于，后者允许在构造过程中使用部分递归函数，只要求最终得到的函数是全函数。因此，“一般”一词（历史遗留的称呼）其实是个误称；从表面上看，\olref{defn:general-recursive} 比 \olref[par]{defn:recursive-fn} \emph{更不}一般。但幸运的是，这种差别只是假象；尽管定义不同，一般递归函数集合与递归函数集合实质上是同一个集合。

\end{document}